\documentclass[10pt,a4paper]{amsart}
\usepackage{amsmath,amssymb,amsthm,graphicx}
% Retained preamble for the Stacks Project chapter snapshot in this directory.
% The chapter's own preamble.tex is a copy of the shared Stacks Project preamble; the block
% below is sliced MECHANICALLY, line for line, from it. Nothing was added, removed or
% redefined; the wrapper class, the xr-hyper/hyperref wiring and the externaldocument list
% are omitted because the assembler supplies the document class.
%
% lines taken from preamble.tex: 3, 5, 9, 11, 12, 156, 157, 158, ... (51 lines total)
%
% package coverage restored from the chapter's own preamble (2026-10-13): xy, xy-2cell

\usepackage[all]{xy}
\xyoption{2cell}
\UseAllTwocells

\newenvironment{reference}{\comment}{\endcomment}
\newenvironment{slogan}{\comment}{\endcomment}
\newenvironment{history}{\comment}{\endcomment}
\theoremstyle{plain}
\newtheorem{theorem}[subsection]{Theorem}
\newtheorem{proposition}[subsection]{Proposition}
\newtheorem{lemma}[subsection]{Lemma}
\theoremstyle{definition}
\newtheorem{definition}[subsection]{Definition}
\newtheorem{example}[subsection]{Example}
\newtheorem{exercise}[subsection]{Exercise}
\newtheorem{situation}[subsection]{Situation}
\theoremstyle{remark}
\newtheorem{remark}[subsection]{Remark}
\newtheorem{remarks}[subsection]{Remarks}
\numberwithin{equation}{subsection}
\def\lim{\mathop{\mathrm{lim}}\nolimits}
\def\colim{\mathop{\mathrm{colim}}\nolimits}
\def\Spec{\mathop{\mathrm{Spec}}}
\def\Hom{\mathop{\mathrm{Hom}}\nolimits}
\def\Ext{\mathop{\mathrm{Ext}}\nolimits}
\def\SheafHom{\mathop{\mathcal{H}\!\mathit{om}}\nolimits}
\def\SheafExt{\mathop{\mathcal{E}\!\mathit{xt}}\nolimits}
\def\Sch{\mathit{Sch}}
\def\Mor{\mathop{\mathrm{Mor}}\nolimits}
\def\Ob{\mathop{\mathrm{Ob}}\nolimits}
\def\Sh{\mathop{\mathit{Sh}}\nolimits}
\def\NL{\mathop{N\!L}\nolimits}
\def\CH{\mathop{\mathrm{CH}}\nolimits}
\def\proetale{{pro\text{-}\acute{e}tale}}
\def\etale{{\acute{e}tale}}
\def\QCoh{\mathit{QCoh}}
\def\Ker{\mathop{\mathrm{Ker}}}
\def\Im{\mathop{\mathrm{Im}}}
\def\Coker{\mathop{\mathrm{Coker}}}
\def\Coim{\mathop{\mathrm{Coim}}}
\DeclareMathSymbol{\boxtimes}{\mathbin}{AMSa}{"02}
\def\QCohstack{\mathcal{QC}\!\mathit{oh}}
\def\Cohstack{\mathcal{C}\!\mathit{oh}}
\def\Spacesstack{\mathcal{S}\!\mathit{paces}}
\def\Quotfunctor{\mathrm{Quot}}
\def\Hilbfunctor{\mathrm{Hilb}}
\def\Curvesstack{\mathcal{C}\!\mathit{urves}}
\def\Polarizedstack{\mathcal{P}\!\mathit{olarized}}
\def\Complexesstack{\mathcal{C}\!\mathit{omplexes}}
\def\Pic{\mathop{\mathrm{Pic}}\nolimits}
\def\Picardstack{\mathcal{P}\!\mathit{ic}}
\def\Picardfunctor{\mathrm{Pic}}
\def\Deformationcategory{\mathcal{D}\!\mathit{ef}}
\usepackage{cleveref}
\crefname{originallemma}{original lemma}{original lemmas}
\Crefname{originallemma}{Original lemma}{Original lemmas}
\makeatletter
\newlabel{duality-lemma-duality-bootstrap}{{0AUE}{source}{}{sourcepointer-1}{}}
\newlabel{duality-lemma-duality-bootstrap@cref}{{[originallemma][0][]0AUE}{[1][0][]source}{}{}{}}
\makeatother
\crefname{originallemma}{original lemma}{original lemmas}
\Crefname{originallemma}{Original lemma}{Original lemmas}
\makeatletter
\newlabel{duality-lemma-finite-twisted}{{0AX2}{source}{}{sourcepointer-2}{}}
\newlabel{duality-lemma-finite-twisted@cref}{{[originallemma][0][]0AX2}{[1][0][]source}{}{}{}}
\makeatother
\crefname{originallemma}{original lemma}{original lemmas}
\Crefname{originallemma}{Original lemma}{Original lemmas}
\makeatletter
\newlabel{duality-lemma-shriek-well-defined}{{0AA0}{source}{}{sourcepointer-3}{}}
\newlabel{duality-lemma-shriek-well-defined@cref}{{[originallemma][0][]0AA0}{[1][0][]source}{}{}{}}
\makeatother
\crefname{originallemma}{original lemma}{original lemmas}
\Crefname{originallemma}{Original lemma}{Original lemmas}
\makeatletter
\newlabel{duality-lemma-twisted-inverse-image}{{0A9E}{source}{}{sourcepointer-4}{}}
\newlabel{duality-lemma-twisted-inverse-image@cref}{{[originallemma][0][]0A9E}{[1][0][]source}{}{}{}}
\makeatother
\crefname{originallemma}{original lemma}{original lemmas}
\Crefname{originallemma}{Original lemma}{Original lemmas}
\makeatletter
\newlabel{duality-section-upper-shriek}{{0A9Y}{source}{}{sourcepointer-5}{}}
\newlabel{duality-section-upper-shriek@cref}{{[originallemma][0][]0A9Y}{[1][0][]source}{}{}{}}
\makeatother
\crefname{originallemma}{original lemma}{original lemmas}
\Crefname{originallemma}{Original lemma}{Original lemmas}
\makeatletter
\newlabel{lemma-localize-dualizing}{{0BT2}{source}{}{sourcepointer-6}{}}
\newlabel{lemma-localize-dualizing@cref}{{[originallemma][0][]0BT2}{[1][0][]source}{}{}{}}
\makeatother
\crefname{originallemma}{original lemma}{original lemmas}
\Crefname{originallemma}{Original lemma}{Original lemmas}
\makeatletter
\newlabel{more-morphisms-lemma-quasi-finite-separated-pass-through-finite}{{05K0}{source}{}{sourcepointer-7}{}}
\newlabel{more-morphisms-lemma-quasi-finite-separated-pass-through-finite@cref}{{[originallemma][0][]05K0}{[1][0][]source}{}{}{}}
\makeatother
\crefname{originallemma}{original lemma}{original lemmas}
\Crefname{originallemma}{Original lemma}{Original lemmas}
\makeatletter
\newlabel{morphisms-lemma-affine-equivalence-modules}{{01SB}{source}{}{sourcepointer-8}{}}
\newlabel{morphisms-lemma-affine-equivalence-modules@cref}{{[originallemma][0][]01SB}{[1][0][]source}{}{}{}}
\makeatother
\crefname{originallemma}{original lemma}{original lemmas}
\Crefname{originallemma}{Original lemma}{Original lemmas}
\makeatletter
\newlabel{remark-relative-dualizing-for-quasi-finite}{{0BVG}{source}{}{sourcepointer-9}{}}
\newlabel{remark-relative-dualizing-for-quasi-finite@cref}{{[originallemma][0][]0BVG}{[1][0][]source}{}{}{}}
\makeatother
\title{Stacks, Discriminants: lemma-compare-dualizing}
\author{The Stacks Project Authors}
\date{stacks-project:discriminant.tex:local-snapshot-20261012; GNU FDL 1.2 or later}
\begin{document}
\maketitle
\noindent\textit{Editorial note.} Statement span is lines 3359-3375 of sources/stacks-discriminant/pinned-original-source.tex; proof span is lines 3377-3431, the complete author proof (53 body lines).\par\medskip
\noindent\textit{Editorial note.} Navigation stubs: 9 cited-but-not-included labels are supplied as non-load-bearing pointers with their pinned tags where the snapshot tag index has them.\par\medskip
\noindent\textit{Editorial note.} extra\_packages is empty on purpose: the retained preamble already loads xy with options and a second declaration clashes in the assembler wrapper.\par\medskip
\noindent\textit{Editorial note.} Text is retained verbatim, including the author's own typos and placeholders.\par\medskip
\noindent\textit{Editorial note.} Queue index 5 of queue\_discriminant.json (pb=2373, stmt 3359-3375, proof 3377-3431).\par\medskip
\noindent\textit{Allowed standard prerequisites (editorial).}\begin{itemize}
\item Explicit prerequisite boundary. The cited inputs are prerequisites; the argument in the retained span is not replaced by any cited result.
\end{itemize}\medskip
\ifdefined\hypertarget\hypertarget{sourcepointer-1}{}\fi
\noindent\textit{Source reference.} Original lemma 0AUE: sources/stacks-discriminant/pinned-original-source.tex, source lines 3431--3431. Cited by the retained statement or proof but not part of the collected spans; supplied only as a navigation pointer. The argument does not depend on its content.\par\medskip
\ifdefined\hypertarget\hypertarget{sourcepointer-2}{}\fi
\noindent\textit{Source reference.} Original lemma 0AX2: sources/stacks-discriminant/pinned-original-source.tex, source lines 3431--3431. Cited by the retained statement or proof but not part of the collected spans; supplied only as a navigation pointer. The argument does not depend on its content.\par\medskip
\ifdefined\hypertarget\hypertarget{sourcepointer-3}{}\fi
\noindent\textit{Source reference.} Original lemma 0AA0: sources/stacks-discriminant/pinned-original-source.tex, source lines 3431--3431. Cited by the retained statement or proof but not part of the collected spans; supplied only as a navigation pointer. The argument does not depend on its content.\par\medskip
\ifdefined\hypertarget\hypertarget{sourcepointer-4}{}\fi
\noindent\textit{Source reference.} Original lemma 0A9E: sources/stacks-discriminant/pinned-original-source.tex, source lines 3431--3431. Cited by the retained statement or proof but not part of the collected spans; supplied only as a navigation pointer. The argument does not depend on its content.\par\medskip
\ifdefined\hypertarget\hypertarget{sourcepointer-5}{}\fi
\noindent\textit{Source reference.} Original lemma 0A9Y: sources/stacks-discriminant/pinned-original-source.tex, source lines 3431--3431. Cited by the retained statement or proof but not part of the collected spans; supplied only as a navigation pointer. The argument does not depend on its content.\par\medskip
\ifdefined\hypertarget\hypertarget{sourcepointer-6}{}\fi
\noindent\textit{Source reference.} Original lemma 0BT2: sources/stacks-discriminant/pinned-original-source.tex, source lines 380--387. Cited by the retained statement or proof but not part of the collected spans; supplied only as a navigation pointer. The argument does not depend on its content.\par\medskip
\ifdefined\hypertarget\hypertarget{sourcepointer-7}{}\fi
\noindent\textit{Source reference.} Original lemma 05K0: sources/stacks-discriminant/pinned-original-source.tex, source lines 3431--3431. Cited by the retained statement or proof but not part of the collected spans; supplied only as a navigation pointer. The argument does not depend on its content.\par\medskip
\ifdefined\hypertarget\hypertarget{sourcepointer-8}{}\fi
\noindent\textit{Source reference.} Original lemma 01SB: sources/stacks-discriminant/pinned-original-source.tex, source lines 3431--3431. Cited by the retained statement or proof but not part of the collected spans; supplied only as a navigation pointer. The argument does not depend on its content.\par\medskip
\ifdefined\hypertarget\hypertarget{sourcepointer-9}{}\fi
\noindent\textit{Source reference.} Original lemma 0BVG: sources/stacks-discriminant/pinned-original-source.tex, source lines 649--660. Cited by the retained statement or proof but not part of the collected spans; supplied only as a navigation pointer. The argument does not depend on its content.\par\medskip
\begin{lemma}
\label{lemma-compare-dualizing}
Let $f : Y \to X$ be a quasi-finite separated morphism of Noetherian schemes.
For every pair of affine opens $\Spec(B) = V \subset Y$,
$\Spec(A) = U \subset X$ with $f(V) \subset U$ there is an isomorphism
$$
H^0(V, f^!\mathcal{O}_X) = \omega_{B/A}
$$
where $f^!$ is as in
Duality for Schemes, Section \ref{duality-section-upper-shriek}.
These isomorphisms are compatible with restriction maps and define a canonical
isomorphism $H^0(f^!\mathcal{O}_X) = \omega_{Y/X}$ with
$\omega_{Y/X}$ as in Remark \ref{remark-relative-dualizing-for-quasi-finite}.
Similarly, if $f : Y \to X$ is a quasi-finite morphism of schemes of
finite type over a Noetherian base $S$ endowed with a dualizing complex
$\omega_S^\bullet$, then $H^0(f_{new}^!\mathcal{O}_X) = \omega_{Y/X}$.
\end{lemma}

\begin{proof}
By Zariski's main theorem we can choose a factorization $f = f' \circ j$
where $j : Y \to Y'$ is an open immersion and $f' : Y' \to X$ is a finite
morphism, see More on Morphisms, Lemma
\ref{more-morphisms-lemma-quasi-finite-separated-pass-through-finite}.
By our construction in
Duality for Schemes, Lemma \ref{duality-lemma-shriek-well-defined} we have
$f^! = j^* \circ a'$ where
$a' : D_\QCoh(\mathcal{O}_X) \to D_\QCoh(\mathcal{O}_{Y'})$
is the right adjoint to $Rf'_*$ of
Duality for Schemes, Lemma \ref{duality-lemma-twisted-inverse-image}.
By Duality for Schemes, Lemma \ref{duality-lemma-finite-twisted}
we see that
$\Phi(a'(\mathcal{O}_X)) = R\SheafHom(f'_*\mathcal{O}_{Y'}, \mathcal{O}_X)$ in
$D_\QCoh^+(f'_*\mathcal{O}_{Y'})$. In particular $a'(\mathcal{O}_X)$ has
vanishing cohomology sheaves in degrees $< 0$. The zeroth cohomology sheaf
is determined by the isomorphism
$$
f'_*H^0(a'(\mathcal{O}_X)) =
\SheafHom_{\mathcal{O}_X}(f'_*\mathcal{O}_{Y'}, \mathcal{O}_X)
$$
as $f'_*\mathcal{O}_{Y'}$-modules via the equivalence of
Morphisms, Lemma \ref{morphisms-lemma-affine-equivalence-modules}.
Writing $(f')^{-1}U = V' = \Spec(B')$, we obtain
$$
H^0(V', a'(\mathcal{O}_X)) = \Hom_A(B', A).
$$
As the zeroth cohomology sheaf of $a'(\mathcal{O}_X)$
is a quasi-coherent module we find that
the restriction to $V$ is given by
$\omega_{B/A} = \Hom_A(B', A) \otimes_{B'} B$ as desired.

\medskip\noindent
The statement about restriction maps signifies that the restriction mappings
of the quasi-coherent $\mathcal{O}_{Y'}$-module $H^0(a'(\mathcal{O}_X))$
for opens in $Y'$ agrees with the maps defined in
Lemma \ref{lemma-localize-dualizing}
for the modules $\omega_{B/A}$ via the isomorphisms given above.
This is clear.

\medskip\noindent
Let $f : Y \to X$ be a quasi-finite morphism of schemes of finite type
over a Noetherian base $S$ endowed with a dualizing complex $\omega_S^\bullet$.
Consider opens $V \subset Y$ and $U \subset X$ with $f(V) \subset U$
and $V$ and $U$ separated over $S$. Denote $f|_V : V \to U$ the restriction
of $f$. By the discussion above and
Duality for Schemes, Lemma \ref{duality-lemma-duality-bootstrap}
there are canonical isomorphisms
$$
H^0(f_{new}^!\mathcal{O}_X)|_V = H^0((f|_V)^!\mathcal{O}_U) = \omega_{V/U} =
\omega_{Y/X}|_V
$$
We omit the verification that these isomorphisms glue to a global
isomorphism $H^0(f_{new}^!\mathcal{O}_X) \to \omega_{Y/X}$.
\end{proof}

\end{document}
